\documentclass[11pt]{article}
\usepackage[margin=1in]{geometry}
\usepackage{amsmath,amssymb,amsthm,booktabs,microtype}
\usepackage[hidelinks]{hyperref}
\newtheorem{theorem}{Theorem}[section]
\newtheorem{proposition}[theorem]{Proposition}
\newtheorem{corollary}[theorem]{Corollary}
\newtheorem{remark}[theorem]{Remark}
\title{Adaptive Stochastic Domination for Yang--Mills Bad Blocks\\
\large The exact quantifier missing from connected-event envelopes}
\author{}
\date{20 July 2026}
\begin{document}
\maketitle

\begin{abstract}
We isolate and prove the probabilistic conversion theorem required by a
dependent Yang--Mills bad-block programme.  If an ordered Boolean bad field
has next-site conditional probability bounded by $q_i$ after every
positive-probability exposure history, then a shared-uniform sequential
coupling places it below independent Bernoulli variables of parameters $q_i$.
This yields joint and connected-cluster bounds without constructing
Mayer/Ursell activities.  A two-site anticorrelated law proves that bounds on
all intersection moments alone do not imply stochastic domination.  At the
Stage 14 effective weight, the abstract degree-$18$ cluster gate is
subcritical.  The remaining task is the full nonlinear $\mathrm{SU}(2)$
adaptive-history estimate.
\end{abstract}

\section{Sequential coupling theorem}

Let $X=(X_1,\ldots,X_N)\in\{0,1\}^N$ have an arbitrary law.  Fix an exposure
order.  For every history $h=(x_1,\ldots,x_{i-1})$ of positive probability,
write
\[
 p_i(h)=\mathbb P(X_i=1\mid X_1=x_1,\ldots,X_{i-1}=x_{i-1}).
\]

\begin{theorem}[Adaptive Bernoulli domination]
If $p_i(h)\leq q_i$ for every $i$ and every positive-probability history,
then there is a coupling with independent
$Y_i\sim\operatorname{Bernoulli}(q_i)$ such that
\[
 X_i\leq Y_i\quad\text{for every }i\quad\text{almost surely}.
\]
Consequently the law of $X$ is stochastically dominated by the product law of
$Y$.
\end{theorem}
\begin{proof}
Take independent $U_i\sim\operatorname{Uniform}[0,1]$.  Generate $X$
sequentially by
\[
 X_i=\mathbf 1\{U_i\leq p_i(X_1,\ldots,X_{i-1})\}.
\]
This reproduces the prescribed conditional probabilities and hence the law of
$X$.  On the same space put $Y_i=\mathbf 1\{U_i\leq q_i\}$.  The $Y_i$ are
independent, and $p_i\leq q_i$ gives $X_i\leq Y_i$ pointwise.
\end{proof}

The premise is history-uniform, not merely an unconditional one-site or joint
moment bound.  It is also compatible with adaptive block selection: the next
site may be a measurable function of the exposed history, provided the same
conditional ceiling holds.

\section{Connected clusters}

Assume $q_i\leq q$ and the bad-block adjacency graph has maximum degree
$\Delta$.  The coupling gives, for every finite set $S$,
\[
 \mathbb P(X_i=1\ \forall i\in S)\leq q^{|S|}.
\]
Using the declared rooted connected-set count
$\Delta^{2(n-1)}$, one obtains
\begin{corollary}
If $e^\alpha\Delta^2q<1$, the weighted sum of bad connected sets containing a
fixed root is at most
\begin{equation}
 \frac{e^\alpha q}{1-e^\alpha\Delta^2q}.
 \label{eq:cluster}
\end{equation}
\end{corollary}

This is a stochastic-domination route.  No partition-function activity is
identified or required for the support-cluster statement.

\section{Why intersection moments are insufficient}

\begin{proposition}[Two-site counterexample]
Let
\[
 \mathbb P(X=(1,0))=\mathbb P(X=(0,1))=\frac12.
\]
For $q=1/2$, every nonempty intersection satisfies
$\mathbb P(X_i=1\ \forall i\in S)\leq q^{|S|}$.  Nevertheless the law is not
stochastically dominated by independent Bernoulli$(1/2)$ variables.
\end{proposition}
\begin{proof}
The singleton probabilities equal $1/2$ and the double intersection has
probability zero.  But the increasing event $\{X_1=1\text{ or }X_2=1\}$ has
probability one, whereas its product-law probability is $3/4$.  Equivalently,
$\mathbb P(X_2=1\mid X_1=0)=1$.
\end{proof}

Thus the connected-event envelopes from earlier stages do not alone justify
stochastic domination.  The adaptive conditional quantifier is essential.

\section{Yang--Mills budget and frontier}

At the Stage 14 sufficient parameters,
\[
 q_{\rm eff}=e^{\kappa-J}+e^{-K}=0.00243027750184,
\]
while for $\alpha=0.1$ and $\Delta=18$,
\[
 q_{\rm crit}=e^{-\alpha}\Delta^{-2}=0.00279270808036.
\]
The branching ratio is $0.870222534$, and \eqref{eq:cluster} is
$0.0206959813$.  Therefore the event-to-cluster conversion closes exactly
\emph{if} every adaptive $\mathrm{SU}(2)$ block history obeys the stated
conditional bound.

That premise remains open.  The corrected frontier is to construct a
gauge-covariant multi-link exposure retaining sufficient boundary information
to charge center and continuous collar frustration, and to prove
\[
 \mathbb P(\text{next block bad}\mid\text{exposed history})
 \leq0.00243027750184
\]
uniformly in volume.  Repeated RG, continuum construction, and the mass gap
remain downstream.  The present theorem is C2.

\end{document}
